\chapter{面向物理仿真的三维场景构建与材质初始化}

\section{引言}
构建能够支撑物理仿真的高保真三维场景，是本文系统的基础模块。传统的人工虚拟环境建模，效率低，成本高。随着生成式 AI 的发展，基于单图或文本自动构建三维几何场景成为可能。然而，直接构建的三维分布仅具备视觉表象，缺乏物理参数与本构分类，无法支持后续物理求解器的直接介入。
为解决上述问题，本章实现了物理参数赋予的流程。FlashWorld 场景生成架构完成初始场景的生成，利用依托分割大模型（SAM）与视觉语言大模型（Qwen3-VL-32B\cite{qwen3vl}），实现对场景多视角图像的分割和物理材质预测，并结合材质先验库，辅助推理出各个材质区域的先验杨氏模量、泊松比与密度初始解。流程图如图~\ref{fig:vlm}~所示。

\begin{figure}[!htbp]
    \centering
    \includegraphics[width=\linewidth]{sam_vlm.png}
    \caption{基于视觉语言模型的材质识别流程}
    \label{fig:vlm}
\end{figure}


\section{物理仿真导向的三维场景生成}

\subsection{场景生成}
为了克服人工建模与隐式辐射场在训练和渲染速度上的局限，系统以 FlashWorld\cite{flashworld} 架构实现文本到三维场景的生成。FlashWorld 需要相机序列作为输入，本文采用 Director3D\cite{director3d} 的 Traj\_dit 模块生成多视角观测序列并辅以稳定平滑后处理。Traj\_dit 是一个基于 DiT（Diffusion Transformer）架构的条件扩散模型，以文本描述或参考图像为条件，通过迭代去噪过程生成一组包含相机外参（位置与朝向）的连续轨迹序列。

Traj\_dit 根据输入提示或参考图像生成一段初始的相机运动轨迹,此时的相机轨迹的视角变化过于剧烈，本文在此基础上实施了一系列后处理优化：（1）对生成轨迹进行三次样条插值，填充帧间间隙并缓解运动突变，确保相邻视角之间的相机位移与旋转角平滑过渡；（2）对所有相机位姿执行朝向归一化与视口覆盖校验，剔除指向天空或纯色墙面等低信息量区域的无效视角，保证最低的视口面积覆盖场景有效几何；（3）对通过筛选的轨迹进行密度自适应重采样，在场景复杂度较高的区域适当加密观测点以提升重建质量。经此后处理管线，系统最终获得约 30–60 帧质量稳定、空间分布均匀的多视角观测序列。

综合以上，在给定文本 Prompt 或单张自然图像的情况下，场景生成模块能合成具备一致性视角的多帧观测序列，并同步输出各视角所需的内参、外参及深度信息。这些多帧影像进一步反投影至三维空间，融合为包含色彩、法向量及坐标要素的初始点云，作为后续材质预测、特征蒸馏和物理仿真的启动基点。

\section{复杂场景的空间深度重切分策略}

从二维平面到三维空间的投射不仅会遇到视角缺失问题，还常常遭遇严重的物体遮挡耦合。当依赖二维图像分割进行类别标定时，距离相机极远的天空或桌面等背景，与位于焦平面的前景（例如被前景物体遮挡的碎石），在屏幕空间上常常不可避免地发生重叠，导致后续材质分类产生混淆。

为了恢复图像像素的空间层级约束，系统集成了基础深度模型 Depth-Anything-V2\cite{depthanythingv2}，系统在获取 FlashWorld 多视角图的同时推理对应的逐像素单目深度图。该深度图具备良好的仿射稳健性与层级平滑性，为空间解耦提供相对可靠的尺度基准。

系统以 SAM 语义块为单元进行深度分层。对于 Depth-Anything-V2 输出并经百分位归一化至 $[0,1]$ 的单目深度图与 SAM 分割产生的语义块集合，取各语义块 $s$ 内深度值的第 75 分位数 $d_s$ 作为其代表深度，以抑制遮挡边界处的离群噪声。对视图内所有 $\{d_s\}$ 执行 K-Means\cite{kmeans} 聚类（$K=3$），依聚类质心深度升序映射为远景（$l=0$）、中景（$l=1$）、近景（$l=2$）三层。各语义块整块归入其所属层，输出至 \texttt{gpt\_input\_layer\_\{0,1,2\}/}$\;$与 \texttt{seg\_layer\_\{0,1,2\}/}，供后续 VLM 逐层材质分类。该聚类以语义块而非像素为原子单位，天然继承 SAM 的轮廓精度，硬分配策略在物体边界处不产生人工断裂，为下游 VLM 材质识别提供了几何纯净、空间隔离的逐层输入。

\section{视觉语言大模型驱动的宏观材质识别}

深度分层为空间部件识别提供了物理隔离保证。随后，系统以自动化方法从这些纯净的空间图层中提取独立物体模块，并标定其物理属性标签。

\subsection{SAM 自动语义分割}
系统采用 SAM 自动掩码生成器（SamAutomaticMaskGenerator）对每帧多视角图像进行预分割。该生成器以均匀网格点采样空间提示（points\_per\_side=32），经预测 IoU 阈值（0.7）与稳定性得分阈值（0.85）筛选，输出场景的候选语义块集合 $\mathcal{S}$。各语义块被编码为独立的掩码区域并裁剪为图块，写入 \texttt{gpt\_input/} 目录。相较于依赖预定义类别的检测器方案，自动网格模式无需物体类别先验，在开放场景中可覆盖更广泛的候选区域，产生的语义块天然遵循视觉边界，为后续深度重切分与 VLM 材质识别提供了轮廓完整的操作单元。

\subsection{基于 Qwen3-VL 的逐层材质分类}
系统以 Qwen3-VL-32B\cite{qwen3vl} 作为材质识别 VLM。对每个语义块，VLM 接收一帧三面板拼接图像：左面板为原始场景上下文，中面板为带半透明掩码与红色轮廓的语义目标区域，右面板为黑背景裁剪图块。提示语要求模型从 9 类预设材质库（石材/陶瓷、玻璃、塑料、金属、木材、橡胶、沙土颗粒、流体、织物）中识别目标区域的主导材质，并以刚性、密度、透明度、反光率、形变倾向等物理线索为决策依据，而非依赖物体身份标签。同时输出 0--1 置信度与硬度等级（soft/medium/hard）。

系统支持多数投票机制以提升鲁棒性：默认以温度 $T=0$ 进行单次确定性推理；当启用投票（num\_votes=3）时，以 $T\in[0.05, 0.30]$ 的递增温度进行多次采样，取众数为最终材质标签。VLM 后端支持 Qwen3-VL-32B 本地部署与云端 API 两种模式。经深度分层后，语义块按所在的远景/中景/近景层分别送入 VLM，各层独立预测，从架构上避免了前后景材质的跨层混淆。

\subsection{物理参数经验映射}
获得各语义块的材质标签与硬度等级后，系统通过 \texttt{material\_config.json}经验库将语义转化为物理参数。该配置库覆盖 9 类材质，每类材质预定义了杨氏模量 $E$、泊松比 $\nu$、密度 $\rho$ 的 min/max/default 经验范围，并关联对应的 MPM 求解器编号（\texttt{solver\_id}）：例如石材类映射为 Fixed Corotated 弹性求解器（solver\_id=0），橡胶类映射为 Maxwell 粘弹性求解器（solver\_id=9），流体类映射为基于压力的弱可压缩求解器（solver\_id=6）。

系统以组合编码 $\texttt{composite\_id}=\texttt{material\_label}\times10+\texttt{grade\_idx}$ 统一类别与硬度信息，其中 material\_label 取值为 1--9，grade\_idx 取值为 0--2（对应 soft/medium/hard）。例如，VLM 将花盆识别为石材/陶瓷（composite\_id=1×10+1=11），配置器赋予 $E\approx10^8$\,Pa、$\nu\approx0.23$ 的刚性初始区间；帽子识别为织物（composite\_id=8×10+1=81），获得 $E\approx3\times10^6$\,Pa 的低杨氏模量区间。这些区间解构成后续 SVD 视频扩散先验引导的梯度寻优（DreamPhysics\cite{dreamphysics}）的强先验约束边界。

\section{本章小结}
本章详述了系统面对单一视觉生成时物理结构缺失的挑战，如何通过多模块集成完成包含物理信息的场景架构工作。针对初态提取，采用 FlashWorld 配合 traj\_dit 扩散轨迹生成器快速产出多视角图像、相机位姿与初始点云，并通过样条插值与外参校验等后处理确保观测序列的平滑性与一致性。针对遮挡伪影耦合问题，引入 Depth-Anything-V2 单目深度估计与基于 K-Means 的三层语义块深度重切分技术，将交叠的前后景物体约束至独立深度层，消除二维掩码的互相渗透。在此基础上，采用 SAM 自动掩码生成器进行均匀网格点采样分割，获取场景的候选语义块集合。随后利用 Qwen3-VL-32B 视觉语言大模型以三次多数投票机制对每个孤立区域进行材质推理，最终通过 \texttt{material\_config.json} 经验映射与 composite\_id 组合编码体系将语义标签转化为具体的杨氏模量、泊松比、密度及 MPM 求解器类型。这套场景预处理机制为后续的特征蒸馏与动力学演算提供了几何纯净、语义明确且物理参数完备的三维初始场。
